\chapter{BX-128: A Pinned 128-bit Register}
\label{ch:bx128}

\section{The board-level picture}

BX-128 is the BX-64 board with the vector file pinned at 128 bits. It does not replace a console. Instruction fetch has its own 32-bit port. Data memory is a 64-bit port. The hart is RV64, Chapter~\ref{ch:rv64}. The vector file is thirty-two registers of 128 bits, 512 bytes. Chapter~\ref{ch:rv128} is the integer base with XLEN 128. This board keeps the RV64 hart and pins the vector register.

A vector load at SEW 64 moves two elements from one register. A vector load at SEW 32 moves four. Lane counts, the two register files, and this pin beside an RV128 hart are Chapter~\ref{ch:fpga}.

\begin{spec}
XLEN = 64. ELEN = 64. VLEN = 128.\\
Thirty-two vector registers. VRF = 512 bytes.\\
Instruction port 32 bits. Data port 64 bits.
\end{spec}

\section{The architecture}

\begin{spec}
Legal SEW = 8, 16, 32, and 64.\\
VLMAX = 16, 8, 4, and 2 at those widths.\\
FP32 and FP64, including fused multiply-add.\\
If the claim is V: every instruction V requires, and the V bit set in \texttt{misa}.\\
If the claim is Zve64d: do not set the V bit.\\
SEW 128 traps.
\end{spec}

\texttt{vsetvli} with \texttt{rs1} equal to \texttt{x0} writes VLMAX into \texttt{rd}. VLEN is not an instruction field. The acceptance filename records 128, which is the length this image was accepted on.

\section{Widening}

Widening multiply and narrowing shift are the BX-64 operations. They are legal here because ELEN is 64. A 128-bit register holds more elements of that SEW than an eight-bit machine does, and the same number the BX-64 floor holds. In the matrix of Section~\ref{sec:blas} the extra elements are a taller column slice. LMUL groups registers. A group is still a vector.

\section{Opcode map for BX-128}

The parcels are the BX-64 parcels. The diagrams below are the two fields this chapter uses. Integer, FP32, FP64, fused multiply-add, mask, and indexed forms are the map in Chapter~\ref{ch:bx64}.

\subsection{Configuration}

Identical to Chapter~\ref{ch:bx8}. \texttt{zimm} selects SEW and LMUL. Bit 31 stays clear. \texttt{vsetvli t0, zero, e32, m1, tu, mu} is the word \texttt{002072D7}. On this machine \texttt{t0} receives 4.

\subsection{Unit-stride load at SEW=64 --- \texttt{vle64.v}}

\begin{center}
\begingroup
\ttfamily\footnotesize
\setlength{\tabcolsep}{0.3em}
\begin{tabular}{*{16}{w{c}{1.05em}}}
31 & 30 & 29 & 28 & 27 & 26 & 25 & 24 & 23 & 22 & 21 & 20 & 19 & 18 & 17 & 16 \\
\hline
\multicolumn{3}{|c|}{N} &
\multicolumn{1}{c|}{E} &
\multicolumn{2}{c|}{P} &
\multicolumn{1}{c|}{M} &
\multicolumn{5}{c|}{L} &
\multicolumn{4}{c|}{S} \\
\hline
\end{tabular}
\vspace{0.35ex}

\begin{tabular}{*{16}{w{c}{1.05em}}}
15 & 14 & 13 & 12 & 11 & 10 & 9 & 8 & 7 & 6 & 5 & 4 & 3 & 2 & 1 & 0 \\
\hline
\multicolumn{3}{|c|}{W} &
\multicolumn{5}{c|}{D} &
\multicolumn{7}{c|}{O} \\
\hline
\end{tabular}
\endgroup
\end{center}

\vspace{1.4ex}
\begin{center}
\begin{tabular}{@{}ll@{}}
N & = nf \\
E & = mew \\
P & = mop (00 = unit-stride) \\
M & = vm \\
L & = lumop (00000) \\
S & = rs1 (base) \\
W & = width (111 = 64-bit) \\
D & = vd \\
O & = 0000111 \\
\end{tabular}
\end{center}

EEW 8, 16, and 32 keep the width codes already printed: 000, 101, and 110.

\section{RV64 and COP3 fetch}

\begin{center}
\begin{tabular}{@{}lll@{}}
\toprule
Path & Escape & Address \\
\midrule
RV64 & native OP-V & full 64-bit GPR \\
COP3 host & \texttt{4C000000} & low 64 bits of the GPR \\
\bottomrule
\end{tabular}
\end{center}

The RV64 parcel is little-endian. The COP3 escape is the BX-32 escape, and the parcel endian is the host's. COP2 stays the host's own unit. There is no harness directory for this machine.

\section{Assembler kernels that must run}

Every BX-64 acceptance image is required here. The tables stay in Chapter~\ref{ch:bx64}. This chapter prints the tiles whose immediate stride is VLMAX at SEW 32. Offsets 16, 32, and 48 are that stride. \texttt{gemm64} is two doubles.

\begin{table}[htbp]
\centering
\caption{Kernel analogues on this machine.}
\small
\begin{tabular}{@{}llll@{}}
\toprule
Kernel & RVV & Intel & ARM \\
\midrule
gemm32 & \texttt{vfmacc.vf} & \texttt{vbroadcastss} & \texttt{ld1r} \\
gemm64 & \texttt{vfmacc.vf} & \texttt{vmovddup} & \texttt{ld1r} \\
\bottomrule
\end{tabular}
\end{table}

\texttt{gemm32} is the \texttt{sgemm} tile of Chapter~\ref{ch:bx32}: OpenBLAS, BLIS, Eigen, and a convolution after the patches are packed. \texttt{gemm64} is \texttt{dgemm} at two doubles. The rest of that catalogue, including \texttt{wmul32}, \texttt{bfly32}, and \texttt{cmul32}, is required here at the same element width.

\begin{table}[htbp]
\centering
\caption{Kernel comparison: gemm32}
\begingroup
\scriptsize
\setlength{\tabcolsep}{4pt}
\noindent \texttt{a0 = K, a1 = A panel, a2 = B panel, a3 = C. 4-wide outer product.}\par
\vspace{0.6ex}
\begin{tabular}{@{}l|l|l|p{7.2cm}@{}}
\toprule
ISA & Address & Hex & Assembly \\
\midrule
RVV & \texttt{0000} & \texttt{06050A63} & \texttt{gemm32: beqz a0, end} \\
RVV & \texttt{0004} & \texttt{00207057} & \texttt{vsetvli zero, zero, e32, m1, tu, mu} \\
RVV & \texttt{0008} & \texttt{0206E407} & \texttt{vle32.v v8, (a3)} \\
RVV & \texttt{000C} & \texttt{01068293} & \texttt{addi t0, a3, 16} \\
RVV & \texttt{0010} & \texttt{0202E487} & \texttt{vle32.v v9, (t0)} \\
RVV & \texttt{0014} & \texttt{02068293} & \texttt{addi t0, a3, 32} \\
RVV & \texttt{0018} & \texttt{0202E507} & \texttt{vle32.v v10, (t0)} \\
RVV & \texttt{001C} & \texttt{03068293} & \texttt{addi t0, a3, 48} \\
RVV & \texttt{0020} & \texttt{0202E587} & \texttt{vle32.v v11, (t0)} \\
RVV & \texttt{0024} & \texttt{0205E087} & \texttt{loop: vle32.v v1, (a1)} \\
RVV & \texttt{0028} & \texttt{00062007} & \texttt{flw ft0, 0(a2)} \\
RVV & \texttt{002C} & \texttt{B2105457} & \texttt{vfmacc.vf v8, ft0, v1} \\
RVV & \texttt{0030} & \texttt{00462007} & \texttt{flw ft0, 4(a2)} \\
RVV & \texttt{0034} & \texttt{B21054D7} & \texttt{vfmacc.vf v9, ft0, v1} \\
RVV & \texttt{0038} & \texttt{00862007} & \texttt{flw ft0, 8(a2)} \\
RVV & \texttt{003C} & \texttt{B2105557} & \texttt{vfmacc.vf v10, ft0, v1} \\
RVV & \texttt{0040} & \texttt{00C62007} & \texttt{flw ft0, 12(a2)} \\
RVV & \texttt{0044} & \texttt{B21055D7} & \texttt{vfmacc.vf v11, ft0, v1} \\
RVV & \texttt{0048} & \texttt{01058593} & \texttt{addi a1, a1, 16} \\
RVV & \texttt{004C} & \texttt{01060613} & \texttt{addi a2, a2, 16} \\
RVV & \texttt{0050} & \texttt{FFF50513} & \texttt{addi a0, a0, -1} \\
RVV & \texttt{0054} & \texttt{FC0518E3} & \texttt{bnez a0, loop} \\
RVV & \texttt{0058} & \texttt{0206E427} & \texttt{vse32.v v8, (a3)} \\
RVV & \texttt{005C} & \texttt{01068293} & \texttt{addi t0, a3, 16} \\
RVV & \texttt{0060} & \texttt{0202E4A7} & \texttt{vse32.v v9, (t0)} \\
RVV & \texttt{0064} & \texttt{02068293} & \texttt{addi t0, a3, 32} \\
RVV & \texttt{0068} & \texttt{0202E527} & \texttt{vse32.v v10, (t0)} \\
RVV & \texttt{006C} & \texttt{03068293} & \texttt{addi t0, a3, 48} \\
RVV & \texttt{0070} & \texttt{0202E5A7} & \texttt{vse32.v v11, (t0)} \\
RVV & \texttt{0074} & \texttt{00008067} & \texttt{end: ret} \\
\midrule
Intel & \texttt{0000} & \texttt{C5F81006} & \texttt{gemm32: vmovups xmm0, [rsi]} \\
Intel & \texttt{0004} & \texttt{C4E279180A} & \texttt{vbroadcastss xmm1, dword ptr [rdx]} \\
Intel & \texttt{0009} & \texttt{C4E279B8E1} & \texttt{vfmadd231ps xmm4, xmm0, xmm1} \\
Intel & \texttt{000E} & \texttt{C4E279184A04} & \texttt{vbroadcastss xmm1, dword ptr [rdx+4]} \\
Intel & \texttt{0014} & \texttt{C4E279B8E9} & \texttt{vfmadd231ps xmm5, xmm0, xmm1} \\
Intel & \texttt{0019} & \texttt{C4E279184A08} & \texttt{vbroadcastss xmm1, dword ptr [rdx+8]} \\
Intel & \texttt{001F} & \texttt{C4E279B8F1} & \texttt{vfmadd231ps xmm6, xmm0, xmm1} \\
Intel & \texttt{0024} & \texttt{C4E279184A0C} & \texttt{vbroadcastss xmm1, dword ptr [rdx+12]} \\
Intel & \texttt{002A} & \texttt{C4E279B8F9} & \texttt{vfmadd231ps xmm7, xmm0, xmm1} \\
Intel & \texttt{002F} & \texttt{4883C610} & \texttt{add rsi, 16} \\
Intel & \texttt{0033} & \texttt{4883C210} & \texttt{add rdx, 16} \\
Intel & \texttt{0037} & \texttt{4883E901} & \texttt{sub rcx, 1} \\
Intel & \texttt{003B} & \texttt{75C3} & \texttt{jnz gemm32} \\
\midrule
ARM & \texttt{0000} & \texttt{4CDF7820} & \texttt{gemm32: ld1 \{v0.4s\}, [x1], \#16} \\
ARM & \texttt{0004} & \texttt{4DDFC841} & \texttt{ld1r \{v1.4s\}, [x2], \#4} \\
ARM & \texttt{0008} & \texttt{4E21CC04} & \texttt{fmla v4.4s, v0.4s, v1.4s} \\
ARM & \texttt{000C} & \texttt{4DDFC841} & \texttt{ld1r \{v1.4s\}, [x2], \#4} \\
ARM & \texttt{0010} & \texttt{4E21CC05} & \texttt{fmla v5.4s, v0.4s, v1.4s} \\
ARM & \texttt{0014} & \texttt{4DDFC841} & \texttt{ld1r \{v1.4s\}, [x2], \#4} \\
ARM & \texttt{0018} & \texttt{4E21CC06} & \texttt{fmla v6.4s, v0.4s, v1.4s} \\
ARM & \texttt{001C} & \texttt{4DDFC841} & \texttt{ld1r \{v1.4s\}, [x2], \#4} \\
ARM & \texttt{0020} & \texttt{4E21CC07} & \texttt{fmla v7.4s, v0.4s, v1.4s} \\
ARM & \texttt{0024} & \texttt{F1000400} & \texttt{subs x0, x0, \#1} \\
ARM & \texttt{0028} & \texttt{54FFFEC1} & \texttt{b.ne gemm32} \\
\bottomrule
\end{tabular}
\endgroup
\end{table}
\noindent\path{kernels/bx32/gemm32.sew32.lmul1.tu.mu.vlen128.xlen32.le.bin}\par

\begin{table}[htbp]
\centering
\caption{Kernel comparison: gemm64}
\begingroup
\scriptsize
\setlength{\tabcolsep}{4pt}
\noindent \texttt{a0 = K, a1 = A panel, a2 = B panel, a3 = C. 2-wide outer product.}\par
\vspace{0.6ex}
\begin{tabular}{@{}l|l|l|p{7.2cm}@{}}
\toprule
ISA & Address & Hex & Assembly \\
\midrule
RVV & \texttt{0000} & \texttt{06050A63} & \texttt{gemm64: beqz a0, end} \\
RVV & \texttt{0004} & \texttt{00307057} & \texttt{vsetvli zero, zero, e64, m1, tu, mu} \\
RVV & \texttt{0008} & \texttt{0206F407} & \texttt{vle64.v v8, (a3)} \\
RVV & \texttt{000C} & \texttt{01068293} & \texttt{addi t0, a3, 16} \\
RVV & \texttt{0010} & \texttt{0202F487} & \texttt{vle64.v v9, (t0)} \\
RVV & \texttt{0014} & \texttt{02068293} & \texttt{addi t0, a3, 32} \\
RVV & \texttt{0018} & \texttt{0202F507} & \texttt{vle64.v v10, (t0)} \\
RVV & \texttt{001C} & \texttt{03068293} & \texttt{addi t0, a3, 48} \\
RVV & \texttt{0020} & \texttt{0202F587} & \texttt{vle64.v v11, (t0)} \\
RVV & \texttt{0024} & \texttt{0205F087} & \texttt{loop: vle64.v v1, (a1)} \\
RVV & \texttt{0028} & \texttt{00063007} & \texttt{fld ft0, 0(a2)} \\
RVV & \texttt{002C} & \texttt{B2105457} & \texttt{vfmacc.vf v8, ft0, v1} \\
RVV & \texttt{0030} & \texttt{00863007} & \texttt{fld ft0, 8(a2)} \\
RVV & \texttt{0034} & \texttt{B21054D7} & \texttt{vfmacc.vf v9, ft0, v1} \\
RVV & \texttt{0038} & \texttt{01063007} & \texttt{fld ft0, 16(a2)} \\
RVV & \texttt{003C} & \texttt{B2105557} & \texttt{vfmacc.vf v10, ft0, v1} \\
RVV & \texttt{0040} & \texttt{01863007} & \texttt{fld ft0, 24(a2)} \\
RVV & \texttt{0044} & \texttt{B21055D7} & \texttt{vfmacc.vf v11, ft0, v1} \\
RVV & \texttt{0048} & \texttt{01058593} & \texttt{addi a1, a1, 16} \\
RVV & \texttt{004C} & \texttt{02060613} & \texttt{addi a2, a2, 32} \\
RVV & \texttt{0050} & \texttt{FFF50513} & \texttt{addi a0, a0, -1} \\
RVV & \texttt{0054} & \texttt{FC0518E3} & \texttt{bnez a0, loop} \\
RVV & \texttt{0058} & \texttt{0206F427} & \texttt{vse64.v v8, (a3)} \\
RVV & \texttt{005C} & \texttt{01068293} & \texttt{addi t0, a3, 16} \\
RVV & \texttt{0060} & \texttt{0202F4A7} & \texttt{vse64.v v9, (t0)} \\
RVV & \texttt{0064} & \texttt{02068293} & \texttt{addi t0, a3, 32} \\
RVV & \texttt{0068} & \texttt{0202F527} & \texttt{vse64.v v10, (t0)} \\
RVV & \texttt{006C} & \texttt{03068293} & \texttt{addi t0, a3, 48} \\
RVV & \texttt{0070} & \texttt{0202F5A7} & \texttt{vse64.v v11, (t0)} \\
RVV & \texttt{0074} & \texttt{00008067} & \texttt{end: ret} \\
\midrule
Intel & \texttt{0000} & \texttt{C5F91006} & \texttt{gemm64: vmovupd xmm0, [rsi]} \\
Intel & \texttt{0004} & \texttt{C5FB120A} & \texttt{vmovddup xmm1, qword ptr [rdx]} \\
Intel & \texttt{0008} & \texttt{C4E2F9B8E1} & \texttt{vfmadd231pd xmm4, xmm0, xmm1} \\
Intel & \texttt{000D} & \texttt{C5FB124A08} & \texttt{vmovddup xmm1, qword ptr [rdx+8]} \\
Intel & \texttt{0012} & \texttt{C4E2F9B8E9} & \texttt{vfmadd231pd xmm5, xmm0, xmm1} \\
Intel & \texttt{0017} & \texttt{C5FB124A10} & \texttt{vmovddup xmm1, qword ptr [rdx+16]} \\
Intel & \texttt{001C} & \texttt{C4E2F9B8F1} & \texttt{vfmadd231pd xmm6, xmm0, xmm1} \\
Intel & \texttt{0021} & \texttt{C5FB124A18} & \texttt{vmovddup xmm1, qword ptr [rdx+24]} \\
Intel & \texttt{0026} & \texttt{C4E2F9B8F9} & \texttt{vfmadd231pd xmm7, xmm0, xmm1} \\
Intel & \texttt{002B} & \texttt{4883C610} & \texttt{add rsi, 16} \\
Intel & \texttt{002F} & \texttt{4883C220} & \texttt{add rdx, 32} \\
Intel & \texttt{0033} & \texttt{4883E901} & \texttt{sub rcx, 1} \\
Intel & \texttt{0037} & \texttt{75C7} & \texttt{jnz gemm64} \\
\midrule
ARM & \texttt{0000} & \texttt{4CDF7C20} & \texttt{gemm64: ld1 \{v0.2d\}, [x1], \#16} \\
ARM & \texttt{0004} & \texttt{4DDFCC41} & \texttt{ld1r \{v1.2d\}, [x2], \#8} \\
ARM & \texttt{0008} & \texttt{4E61CC04} & \texttt{fmla v4.2d, v0.2d, v1.2d} \\
ARM & \texttt{000C} & \texttt{4DDFCC41} & \texttt{ld1r \{v1.2d\}, [x2], \#8} \\
ARM & \texttt{0010} & \texttt{4E61CC05} & \texttt{fmla v5.2d, v0.2d, v1.2d} \\
ARM & \texttt{0014} & \texttt{4DDFCC41} & \texttt{ld1r \{v1.2d\}, [x2], \#8} \\
ARM & \texttt{0018} & \texttt{4E61CC06} & \texttt{fmla v6.2d, v0.2d, v1.2d} \\
ARM & \texttt{001C} & \texttt{4DDFCC41} & \texttt{ld1r \{v1.2d\}, [x2], \#8} \\
ARM & \texttt{0020} & \texttt{4E61CC07} & \texttt{fmla v7.2d, v0.2d, v1.2d} \\
ARM & \texttt{0024} & \texttt{F1000400} & \texttt{subs x0, x0, \#1} \\
ARM & \texttt{0028} & \texttt{54FFFEC1} & \texttt{b.ne gemm64} \\
\bottomrule
\end{tabular}
\endgroup
\end{table}
\noindent\path{kernels/bx64/gemm64.sew64.lmul1.tu.mu.vlen128.xlen64.le.bin}\par


\section{Harts}

One RV64 hart is the default. An optional second hart is the BX-64 organisation: one hart for the short work, one for a longer kernel, a mapped local memory, and a DMA engine programmed by stores. Both harts see VLEN 128. The DMA does not fetch OP-V. Xsvm stays off on a V-claiming image.
